% R583: EGA I, 0, 2.2.1--2.2.5.
% Traducción: OpenAI Codex, GPT-6.1 Sol, esfuerzo Ultra.
\oldpage[0\textsubscript{I}]{23}

\subsection{Espacios noetherianos}
\label{subsection:0.2.2-fr}

\begin{env}[2.2.1]
\label{0.2.2.1-fr}
Se dice que un espacio topológico $X$ es \emph{noetheriano} si el conjunto de los
abiertos de $X$ satisface la condición \emph{maximal}, o, lo que equivale a lo
mismo, si el conjunto de los cerrados de $X$ satisface la condición \emph{minimal}. Se
dice que $X$ es \emph{localmente noetheriano} si todo $x\in X$ admite una
vecindad que es un subespacio noetheriano.
\end{env}

\begin{env}[2.2.2]
\label{0.2.2.2-fr}
Sea $E$ un conjunto ordenado que satisface la condición \emph{minimal} y sea
$\boldsymbol{P}$ una propiedad de los elementos de $E$ sujeta a la condición siguiente: si
$a\in E$ es tal que, para todo $x<a$, $\boldsymbol{P}(x)$ es cierta, entonces $\boldsymbol{P}(a)$ es
cierta. En estas condiciones, \emph{$\boldsymbol{P}(x)$ es cierta para todo $x\in E$}
(«principio de inducción noetheriana»). En efecto, sea $F$ el conjunto de los
$x\in E$ para los que $\boldsymbol{P}(x)$ es falsa; si $F$ no fuera vacío, tendría un
elemento minimal $a$ y, como entonces $\boldsymbol{P}(x)$ es cierta para todo $x<a$, $\boldsymbol{P}(a)$
también sería cierta, lo que es contradictorio.

Aplicaremos este principio, en particular, cuando $E$ sea un \emph{conjunto
de partes cerradas de un espacio noetheriano}.
\end{env}

\begin{env}[2.2.3]
\label{0.2.2.3-fr}
Todo subespacio de un espacio noetheriano es noetheriano. Recíprocamente, todo
espacio topológico que es unión finita de subespacios noetherianos es noetheriano.
\end{env}

\begin{env}[2.2.4]
\label{0.2.2.4-fr}
Todo espacio noetheriano es cuasicompacto; recíprocamente, todo espacio topológico
en el que todo abierto es cuasicompacto es noetheriano.
\end{env}

\begin{env}[2.2.5]
\label{0.2.2.5-fr}
Un espacio noetheriano tiene un número \emph{finito} de componentes
irreducibles, como se ve por inducción noetheriana.
\end{env}
